\chapter{Conclusiones y líneas futuras}

Cuando empecé a planificar este trabajo tenía claro que quería construir algo que funcionara de verdad, no solo un conjunto de conceptos bien explicados sobre papel. Y creo que lo he conseguido: tengo un clúster Kubernetes en marcha, con políticas de seguridad que bloquean configuraciones peligrosas antes de que lleguen a ejecutarse, que corrigen automáticamente lo que pueden corregir, integrado con ArgoCD para que los despliegues sean declarativos y trazables, y con un pipeline que analiza las imágenes en busca de vulnerabilidades antes de publicar nada. Todo conectado, todo funcionando. Eso me da cierta satisfacción.

Lo que más me ha aportado el trabajo no ha sido, en realidad, ninguna de las herramientas concretas. Ha sido entender cómo trabajan de verdad los equipos de DevSecOps: que la seguridad no es una capa que se añade al final, sino algo que tienes que tener presente desde el primer momento. Eso se ve mucho en la teoría, pero hasta que no lo montas tú mismo, con errores y todo, no lo interioriza de verdad. He aprendido haciendo, y muchas cosas las entendí porque me choqué con ellas.

El hallazgo que más me sorprendió, y del que más contento estoy, lo encontré al montar el pipeline de análisis de vulnerabilidades. Tenía dos herramientas trabajando sobre la misma imagen: Trivy y Grype. Trivy la dio por buena. Grype encontró una vulnerabilidad importante en la imagen de Python y bloqueó el pipeline. Dos herramientas, una imagen, resultados distintos. Podría haber bajado el umbral de Grype para que la imagen pasara y haberme quitado el problema de encima, pero no lo hice: fui al origen, actualicé la imagen base a una versión de Python más reciente, y la vulnerabilidad desapareció. Ver eso en la práctica me enseñó más que cualquier artículo sobre defensa en profundidad. Tiene sentido usar varias capas de seguridad, no porque lo diga un libro, sino porque lo vi pasar.

También me resultó interesante descubrir cómo ArgoCD y Kyverno se complementan para mantener el sistema en el estado correcto. Vi cómo Kyverno corregía configuraciones inseguras antes de que llegaran a validación, y cómo ArgoCD detectaba y revertía cualquier cambio manual en el clúster. Hay una lógica elegante en cómo encajan estas dos herramientas. Y también aprendí que si Kyverno falla o se queda sin recursos, bloquea todo el clúster —porque está diseñado para fallar de forma segura, no para dejar pasar peticiones sin revisar—. Eso lo aprendí de la peor manera: el disco del nodo se llenó, Kyverno se cayó, y el clúster quedó bloqueado. Fue un mal rato. Pero también me enseñó algo que no sale en los manuales: que estas herramientas tienen consecuencias reales cuando algo va mal, y que gestionar recursos es parte del trabajo.

Soy consciente de que hay cosas que se han quedado fuera. No llegué a integrar Sysdig Secure para vigilar lo que pasa en el clúster mientras las aplicaciones están en marcha. La comparativa con OPA Gatekeeper la hice a nivel de código, escribiendo las políticas equivalentes en Rego, pero sin llegar a desplegar Gatekeeper en el clúster, porque tener dos admission controllers activos al mismo tiempo habría sido un lío. Y tampoco implementé la firma de imágenes con Cosign, que habría añadido otra capa de seguridad interesante a la cadena. Son cosas que sé que faltan y que en otro momento abordaría, pero también creo que el trabajo tiene un alcance razonable para lo que es.

De cara al futuro, me quedan varias cosas que me gustaría explorar. Una es integrar de verdad una herramienta de seguridad en tiempo de ejecución en un entorno cloud, que es donde tiene sentido. Otra es automatizar todo el montaje del entorno con Ansible, para poder reproducirlo sin tener que hacer nada a mano, que es como debería funcionar cualquier infraestructura seria. Y también me gustaría montar el mismo pipeline en GitLab CI en lugar de GitHub Actions, porque GitLab es algo que se usa mucho en empresas y me interesa ver las diferencias.

Este trabajo no es un punto final. Es más bien el principio de algo. Sé que me queda mucho por aprender, que hay partes que podría haber hecho mejor, y que el mundo de la seguridad en Kubernetes cambia rápido. Pero salgo con una base que antes no tenía, con la sensación de haber entendido de verdad cómo funcionan estas herramientas, y con ganas de seguir. Y eso, para mí, es suficiente.
